% GENERATED FILE. Edit canonical lessons, then run npm run generate:books.
% canonical-source-hash: fnv1a64:8e961372167550bc
% canonical-lessons: AR-C475-dabu, AR-C475-fahd, AR-C475-qunfudh, AR-C475-sinjab, AR-C475-khuffash

\chapter{Wild animals}
\label{ch:ar-wild-animals}
\hlchaptermodality{475}

\begin{chapteropening}
\textbf{By the end of this chapter:} \emph{I can say a hyena, a cheetah, a hedgehog, a squirrel and a bat.}

Complete the last lesson of chapter 475: I can say a hyena, a cheetah, a hedgehog, a squirrel and a bat.
\end{chapteropening}

\section[\texorpdfstring{\ar{ضبع} (ḍabʿ)}{ضبع (ḍabʿ)}]{\ar{ضبع} (ḍabʿ) — a hyena}

\label{lesson:AR-C475-dabu}



\begin{quote}
Before the new one: say the Arabic for a feather, then the Arabic for to float.
\end{quote}

\subsection*{You'll want to know: \ar{ضبع}}
\textbf{\ar{ضبع}} — \emph{ḍabʿ} — \textquotedblleft{}a hyena\textquotedblright{}.

\begin{quote}
\ar{سمعنا} \ar{صوت} \ar{ضبع} \ar{في} \ar{الليل}.
\end{quote}

\begin{quote}
\emph{samiʿnā ṣawt ḍabʿ fī al-layl.}
\end{quote}

\begin{quote}
We heard a hyena at night.
\end{quote}

\begin{cousinweb}
\textbf{\ar{ضبع}} (\emph{ḍabʿ}) was inherited from Proto-Semitic (EtymDB).
\end{cousinweb}

\begin{grammarlens}[title={What you've built}]
One more word for getting around.
\end{grammarlens}

\subsection*{Your turn}
\begin{itemize}
\raggedright
  \item \emph{Say it:} \emph{ḍabʿ}
  \item \emph{Say it:} \emph{ḍabʿ}, once more
  \item \emph{Say it:} \emph{samiʿnā ṣawt ḍabʿ fī al-layl.} — We heard a hyena at night.
  \item \emph{Recall:} say the Arabic for a feather, then the Arabic for to float, then say \emph{ḍabʿ} again
\end{itemize}

\begin{tcolorbox}[breakable,colback=teal!4,colframe=teal!35!black,title={Before you move on}]
What does \ar{ضبع} mean? (\textquotedblleft{}A hyena\textquotedblright{}.) Say it once more.
\end{tcolorbox}
\section[\texorpdfstring{\ar{فهد} (fahd)}{فهد (fahd)}]{\ar{فهد} (fahd) — a cheetah; a leopard}

\label{lesson:AR-C475-fahd}



\begin{quote}
Before the new one: say the Arabic for to float, then the Arabic for a hyena.
\end{quote}

\subsection*{You'll want to know: \ar{فهد}}
\textbf{\ar{فهد}} — \emph{fahd} — \textquotedblleft{}a cheetah; a leopard\textquotedblright{}.

\begin{quote}
\ar{رأينا} \ar{فهدا} \ar{في} \ar{المحمية}.
\end{quote}

\begin{quote}
\emph{raʾaynā fahdan fī al-maḥmiyya.}
\end{quote}

\begin{quote}
We saw a cheetah in the reserve.
\end{quote}

\begin{cousinweb}
\textbf{\ar{فهد}} (\emph{fahd}) is an old Arabic name; the origin is not settled.
\end{cousinweb}

\begin{grammarlens}[title={What you've built}]
One more word for getting around.
\end{grammarlens}

\subsection*{Your turn}
\begin{itemize}
\raggedright
  \item \emph{Say it:} \emph{fahd}
  \item \emph{Say it:} \emph{fahd}, once more
  \item \emph{Say it:} \emph{raʾaynā fahdan fī al-maḥmiyya.} — We saw a cheetah in the reserve.
  \item \emph{Recall:} say the Arabic for to float, then the Arabic for a hyena, then say \emph{fahd} again
\end{itemize}

\begin{tcolorbox}[breakable,colback=teal!4,colframe=teal!35!black,title={Before you move on}]
What does \ar{فهد} mean? (\textquotedblleft{}A cheetah; a leopard\textquotedblright{}.) Say it once more.
\end{tcolorbox}
\section[\texorpdfstring{\ar{قنفذ} (qunfudh)}{قنفذ (qunfudh)}]{\ar{قنفذ} (qunfudh) — a hedgehog}

\label{lesson:AR-C475-qunfudh}



\begin{quote}
Before the new one: say the Arabic for a hyena, then the Arabic for a cheetah.
\end{quote}

\subsection*{You'll want to know: \ar{قنفذ}}
\textbf{\ar{قنفذ}} — \emph{qunfudh} — \textquotedblleft{}a hedgehog\textquotedblright{}.

\begin{quote}
\ar{وجدنا} \ar{قنفذا} \ar{صغيرا} \ar{في} \ar{الحديقة}.
\end{quote}

\begin{quote}
\emph{wajadnā qunfudhan ṣaghīran fī al-ḥadīqa.}
\end{quote}

\begin{quote}
We found a little hedgehog in the garden.
\end{quote}

\begin{cousinweb}
\textbf{\ar{قنفذ}} (\emph{qunfudh}) was inherited from Proto-Semitic (EtymDB).
\end{cousinweb}

\begin{grammarlens}[title={What you've built}]
One more word for getting around.
\end{grammarlens}

\subsection*{Your turn}
\begin{itemize}
\raggedright
  \item \emph{Say it:} \emph{qunfudh}
  \item \emph{Say it:} \emph{qunfudh}, once more
  \item \emph{Say it:} \emph{wajadnā qunfudhan ṣaghīran fī al-ḥadīqa.} — We found a little hedgehog in the garden.
  \item \emph{Recall:} say the Arabic for a hyena, then the Arabic for a cheetah, then say \emph{qunfudh} again
\end{itemize}

\begin{tcolorbox}[breakable,colback=teal!4,colframe=teal!35!black,title={Before you move on}]
What does \ar{قنفذ} mean? (\textquotedblleft{}A hedgehog\textquotedblright{}.) Say it once more.
\end{tcolorbox}
\section[\texorpdfstring{\ar{سنجاب} (sinjāb)}{سنجاب (sinjāb)}]{\ar{سنجاب} (sinjāb) — a squirrel}

\label{lesson:AR-C475-sinjab}



\begin{quote}
Before the new one: say the Arabic for a cheetah, then the Arabic for a hedgehog.
\end{quote}

\subsection*{You'll want to know: \ar{سنجاب}}
\textbf{\ar{سنجاب}} — \emph{sinjāb} — \textquotedblleft{}a squirrel\textquotedblright{}.

\begin{quote}
\ar{صعد} \ar{سنجاب} \ar{الشجرة} \ar{بسرعة}.
\end{quote}

\begin{quote}
\emph{ṣaʿida sinjāb ash-shajara bi-surʿa.}
\end{quote}

\begin{quote}
A squirrel ran up the tree quickly.
\end{quote}

\begin{cousinweb}
\textbf{\ar{سنجاب}} (\emph{sinjāb}) was borrowed from Persian (EtymDB).
\end{cousinweb}

\begin{grammarlens}[title={What you've built}]
One more word for getting around.
\end{grammarlens}

\subsection*{Your turn}
\begin{itemize}
\raggedright
  \item \emph{Say it:} \emph{sinjāb}
  \item \emph{Say it:} \emph{sinjāb}, once more
  \item \emph{Say it:} \emph{ṣaʿida sinjāb ash-shajara bi-surʿa.} — A squirrel ran up the tree quickly.
  \item \emph{Recall:} say the Arabic for a cheetah, then the Arabic for a hedgehog, then say \emph{sinjāb} again
\end{itemize}

\begin{tcolorbox}[breakable,colback=teal!4,colframe=teal!35!black,title={Before you move on}]
What does \ar{سنجاب} mean? (\textquotedblleft{}A squirrel\textquotedblright{}.) Say it once more.
\end{tcolorbox}
\section[\texorpdfstring{\ar{خفاش} (khuffāsh)}{خفاش (khuffāsh)}]{\ar{خفاش} (khuffāsh) — a bat (the animal)}

\label{lesson:AR-C475-khuffash}



\begin{quote}
Before the new one: say the Arabic for a hedgehog, then the Arabic for a squirrel.
\end{quote}

\subsection*{You'll want to know: \ar{خفاش}}
\textbf{\ar{خفاش}} — \emph{khuffāsh} — \textquotedblleft{}a bat (the animal)\textquotedblright{}.

\begin{quote}
\ar{يطير} \ar{الخفاش} \ar{في} \ar{الليل}.
\end{quote}

\begin{quote}
\emph{yaṭīru al-khuffāsh fī al-layl.}
\end{quote}

\begin{quote}
The bat flies at night.
\end{quote}

\begin{cousinweb}
\textbf{\ar{خفاش}} (\emph{khuffāsh}) is built on the native Arabic root \emph{kh-f-sh}, \textquotedblleft{}to have weak eyes\textquotedblright{} (Lane). EtymDB also derives it from this root: the bat was the animal that sees poorly by day.
\end{cousinweb}

\begin{grammarlens}[title={What you've built}]
One more word for getting around.
\end{grammarlens}

\subsection*{Your turn}
\begin{itemize}
\raggedright
  \item \emph{Say it:} \emph{khuffāsh}
  \item \emph{Say it:} \emph{khuffāsh}, once more
  \item \emph{Say it:} \emph{yaṭīru al-khuffāsh fī al-layl.} — The bat flies at night.
  \item \emph{Recall:} say the Arabic for a hedgehog, then the Arabic for a squirrel, then say \emph{khuffāsh} again
\end{itemize}

\begin{tcolorbox}[breakable,colback=teal!4,colframe=teal!35!black,title={Before you move on}]
What does \ar{خفاش} mean? (\textquotedblleft{}A bat (the animal)\textquotedblright{}.) Say it once more.
\end{tcolorbox}
